\documentclass{article}
\usepackage{graphicx} % Required for inserting images
\usepackage[a4paper, left=1.6cm, right=1.6cm, top=1.8cm, bottom=2cm]{geometry}
\usepackage{tikz}
\usepackage{float}
\usepackage{amsmath}
\usepackage{booktabs}
\usepackage{tabularx}
\usepackage{amssymb}
\usetikzlibrary{positioning, arrows.meta, shapes.geometric}
\title{Kimi k2.5 LLM Report}
\author{Zeina Mostafa Ali }


\begin{document}

\maketitle

\section{Introduction}
Moonshot AI developed Kimi K2.5, released on January 27, 2026. It is particularly interesting because it is not simply a text LLM with an image model attached. Moonshot describes it as a native multimodal agentic model, trained through continued pretraining on approximately 15 trillion mixed visual and text tokens on top of Kimi K2 Base.

The main ideas behind Kimi K2.5 are:

\begin{itemize}
\item Mixture-of-Experts (MoE) architecture
\item Native multimodal understanding
\item Long-context processing
\item Strong reasoning and coding capabilities
\item Support for tool use and agentic workflows
\item Agent Swarm for solving complex tasks using multiple agents
\end{itemize}

\section{Architecture}
Kimi K2.5 uses a Mixture-of-Experts (MoE) architecture. Instead of activating all of its parameters for every input token, the model uses a routing mechanism to select a subset of experts.

The main architectural specifications are:

\begin{table}[h]
\centering
\begin{tabular}{ll}

Property & Kimi K2.5 \\

Total parameters & Approximately 1 trillion \\
Active parameters & Approximately 32 billion per token \\
Number of layers & 61 \\
Attention heads & 64 \\
Number of experts & 384 \\
Selected experts per token & 8 \\
Vocabulary size & 160K tokens \\
Attention mechanism & Multi-head Latent Attention (MLA) \\
Activation function & SwiGLU \\
Vision encoder & MoonViT \\
Context length & 256K tokens \\
\hline
\end{tabular}
\caption{Main architectural characteristics of Kimi K2.5}
\end{table}

The important point is that Kimi K2.5 has approximately 1 trillion total parameters, but only around 32 billion parameters are activated for each token. This makes the model computationally more efficient than a dense model with 1 trillion parameters. However, the model still requires a large amount of memory because the complete set of parameters has to be stored when deploying the model.

\section{Training Approach}
Kimi K2.5 was developed by continuing the pretraining of Kimi K2 Base.

Moonshot AI reports that Kimi K2.5 was trained on approximately 15 trillion mixed visual and text tokens.

This is important because visual information is part of the model's training process rather than being treated as an entirely separate capability.

The training approach therefore allows the model to learn relationships between text and visual information.

This contributes to capabilities such as:

\begin{itemize}
\item Image understanding
\item Document understanding
\item Visual question answering
\item Image-to-code generation
\item Visual reasoning
\item Video understanding
\end{itemize}






\section{Multimodal Capabilities}

Kimi K2.5 is a native multimodal model.

A multimodal model can process more than one type of information. In the case of Kimi K2.5, the main focus is on text and visual information.

For example, a user can provide an image of a user interface and ask the model to explain the interface or generate code for it.

A simplified workflow is:

\[
\text{Image} + \text{Text Prompt}
\rightarrow
\text{Kimi K2.5}
\rightarrow
\text{Understanding or Generated Output}
\]

Some practical multimodal applications include:

\begin{itemize}
\item Understanding images
\item Analyzing documents
\item Understanding screenshots
\item Generating code from visual designs
\item Debugging interfaces using screenshots
\item Video understanding
\end{itemize}

This makes Kimi K2.5 useful for applications where the input is not purely text.

\section{Context Window}

Kimi K2.5 supports a context window of approximately 256K tokens.

A context window determines how much information a model can consider during a single interaction.

A large context window is useful for:

\begin{itemize}
\item Large documents
\item Long conversations
\item Large software repositories
\item Research tasks
\item Long code files
\item Complex agent workflows
\end{itemize}

For example, an AI coding assistant can provide a large amount of source code to the model and ask it to analyze relationships between different parts of the project.

However, a large context window does not guarantee that the model will perfectly understand every piece of information inside it. Retrieval quality and reasoning ability are still important.

\section{Reasoning Capabilities}

Kimi K2.5 includes a thinking mode designed for difficult reasoning tasks.

Moonshot AI reports strong results on several reasoning and mathematical benchmarks.

Some reported results include:

\begin{table}[h]
\centering
\begin{tabular}{ll}
\hline
Benchmark & Kimi K2.5 Score \\
\hline
AIME 2025 & 96.1 \\
GPQA Diamond & 87.6 \\
MMLU-Pro & 87.1 \\
MATH-500 & 98.0 \\
MathVista & 90.1 \\
MathVision & 84.2 \\
\hline
\end{tabular}
\caption{Selected reasoning and mathematical benchmark results}
\end{table}

These benchmarks evaluate different aspects of reasoning.

For example, mathematical benchmarks test mathematical problem solving, while GPQA evaluates difficult scientific questions.

Benchmark results should be interpreted carefully because results can depend on the model configuration, prompting method, tool usage, and evaluation procedure.

\section{Coding Capabilities}

Kimi K2.5 is designed for software development and coding tasks.

It can be used for:

\begin{itemize}
\item Code generation
\item Code explanation
\item Debugging
\item Repository analysis
\item Software engineering agents
\item Coding from screenshots
\item Tool-assisted coding
\end{itemize}

Some reported coding benchmark results are:

\begin{table}[h]
\centering
\begin{tabular}{ll}
\hline
Benchmark & Kimi K2.5 Score \\
\hline
SWE-Bench Verified & 76.8 \\
SWE-Bench Multilingual & 73.0 \\
Terminal-Bench 2.0 & 50.8 \\
LiveCodeBench v6 & 85.0 \\
SciCode & 48.7 \\
\hline
\end{tabular}
\caption{Selected coding benchmark results}
\end{table}

SWE-Bench is especially relevant because it evaluates software engineering tasks using real repositories and issues rather than only asking the model to generate isolated code snippets.

8. Agentic Workflows

One of the most interesting features of Kimi K2.5 is its support for agentic workflows.

An AI agent is a system that can reason about a task, use tools, perform actions, and continue working toward a goal.

A simple LLM workflow can be represented as:

\[
\text{User}
\rightarrow
\text{LLM}
\rightarrow
\text{Answer}
\]

An agentic workflow is more complex:

\[
\text{User}
\rightarrow
\text{Agent}
\rightarrow
\text{Reason}
\rightarrow
\text{Use Tools}
\rightarrow
\text{Observe Results}
\rightarrow
\text{Reason Again}
\rightarrow
\text{Final Answer}
\]

Kimi K2.5 introduces an approach called Agent Swarm.

The main idea is that a complex task can be divided into several smaller tasks that can be handled by specialized agents.

For example:

\[
\begin{array}{c}
\text{Main Agent} \\
\downarrow \\
\begin{array}{ccc}
\text{Research Agent} &
\text{Coding Agent} &
\text{Analysis Agent}
\end{array} \\
\downarrow \\
\text{Main Agent} \\
\downarrow \\
\text{Final Result}
\end{array}
\]

This can be useful for research, coding, browsing, and other tasks that can be divided into independent subtasks.

Moonshot reports improved results on some search benchmarks when Agent Swarm and context management are used.

For example:

\begin{table}[h]
\centering
\begin{tabular}{ll}
\hline
Benchmark & Reported Score \\
\hline
BrowseComp & 60.6 \\
BrowseComp + Context Management & 74.9 \\
BrowseComp + Agent Swarm & 78.4 \\
WideSearch & 72.7 \\
WideSearch + Agent Swarm & 79.0 \\
\hline
\end{tabular}
\caption{Selected agentic benchmark results}
\end{table}

\section{Efficiency}

The MoE architecture is one of the main reasons Kimi K2.5 can have a very large number of total parameters while keeping the number of activated parameters relatively smaller.

For every token, only selected experts are activated.

In simplified form:

\[
\text{Input Token}
\rightarrow
\text{Router}
\rightarrow
\text{Selected Experts}
\rightarrow
\text{Output}
\]

Kimi K2.5 has approximately 1 trillion total parameters but activates approximately 32 billion parameters per token.

This improves computational efficiency compared with activating all 1 trillion parameters for every token.

However, the model is still very large from a deployment perspective.

Therefore, there is a difference between:

\begin{itemize}
\item \textbf{Inference computation}: reduced through sparse expert activation.
\item \textbf{Memory requirement}: still very high because the model contains approximately 1 trillion parameters.
\end{itemize}

\section{Hardware and Deployment}

Kimi K2.5 is significantly larger than models normally deployed on a personal computer.

Moonshot AI provides deployment information for frameworks such as:

\begin{itemize}
\item vLLM
\item SGLang
\item KTransformers
\end{itemize}

The documented vLLM deployment configuration includes multiple high-memory GPUs.

A reported BF16 deployment configuration uses 8 NVIDIA H200 GPUs.

Therefore, running the full model locally is not practical for most students using normal laptops or desktop computers.

For most developers, using an API is considerably easier.

The practical options are:

\[
\text{Kimi K2.5}
\rightarrow
\begin{cases}
\text{Cloud/API} & \text{Practical for most developers}\\
\text{Large GPU Cluster} & \text{Possible for large deployments}\\
\text{Local Deployment} & \text{Requires substantial hardware}
\end{cases}
\]

\section{API Availability}

Kimi K2.5 can be accessed through Moonshot AI's ecosystem and third-party inference platforms.

OpenRouter provides an API interface for accessing different LLMs.

However, model availability on third-party platforms can change.

At the time of this research, the OpenRouter page for Kimi K2.5 indicates that the listed endpoint is going away after September 7, 2026.

Therefore, developers should verify the current API availability before building a production application that depends specifically on Kimi K2.5.

This is also an example of why an application should avoid tightly coupling its code to one model provider.

\section{Licensing}

Kimi K2.5 is released under a Modified MIT License.

The license provides permissions similar to the MIT license but includes an additional attribution requirement for extremely large commercial products or services.

Before using Kimi K2.5 in a commercial application, developers should read the current license carefully and verify that their use case satisfies the license requirements.

\section{Comparison with Other Open or Openly Available Models}

Kimi K2.5 can be compared with other open or openly available models such as GPT-OSS-20B, Qwen3-30B-A3B, and Llama 4 Maverick.

\begin{table}[h]
\centering
\begin{tabular}{llll}
\hline
Model & Size & Context & Multimodal \\
\hline
Kimi K2.5 & $\sim$1T total / 32B active & 256K & Yes \\
GPT-OSS-20B & $\sim$21B / 3.6B active & 131K & No \\
Qwen3-30B-A3B & $\sim$30.5B / 3.3B active & 262K & No \\
Llama 4 Maverick & Large MoE & Up to 1M & Yes \\
\hline
\end{tabular}
\caption{High-level comparison of selected models}
\end{table}

The models have different design goals.

Kimi K2.5 focuses strongly on:

\begin{itemize}
\item Multimodal reasoning
\item Coding
\item Long context
\item Agentic workflows
\end{itemize}

GPT-OSS-20B focuses more on efficient reasoning and coding with a much smaller active parameter count.

Qwen3-30B-A3B provides a relatively efficient MoE architecture with a long context window.

Llama 4 models provide large-scale multimodal capabilities and very large context windows.

Therefore, there is no single model that is best for every application. The appropriate model depends on the requirements of the application, available hardware, cost, latency, and required capabilities.

\section{Practical Applications}

Kimi K2.5 can be used in several practical applications.

\subsection{Coding assistants}

It can help developers:

\begin{itemize}
\item Generate code
\item Debug code
\item Explain code
\item Analyze repositories
\item Solve software engineering issues
\end{itemize}

\subsection{Multimodal applications}

It can be used for:

\begin{itemize}
\item Image analysis
\item Document understanding
\item Screenshot analysis
\item Visual question answering
\item Visual debugging
\end{itemize}

\subsection{Research agents}

Agent Swarm makes the model suitable for research systems where a large problem can be divided into multiple smaller research tasks.

For example:

\[
\text{Research Question}
\rightarrow
\begin{cases}
\text{Agent 1: Find papers}\\
\text{Agent 2: Analyze experiments}\\
\text{Agent 3: Compare results}\\
\text{Agent 4: Summarize evidence}
\end{cases}
\rightarrow
\text{Final Research Report}
\]

\subsection{Long-document analysis}

The large context window makes Kimi K2.5 suitable for applications involving:

\begin{itemize}
\item Large technical documents
\item Software repositories
\item Research papers
\item Long conversations
\item Large collections of related information
\end{itemize}

\section{Current Limitations}

Although Kimi K2.5 is a powerful model, it has several limitations.

First, its approximately 1 trillion total parameters make local deployment difficult and expensive.

Second, benchmark results should not be interpreted as a perfect measure of real-world performance. Results depend on the benchmark, prompting strategy, tools, sampling configuration, and evaluation procedure.

Third, agentic workflows can increase latency and token usage because solving a task may require multiple reasoning steps or multiple agents.

Fourth, a large context window does not guarantee perfect reasoning over all provided information.

Finally, API availability can change. Developers should check the current provider and model availability before depending on Kimi K2.5 in production.

\section{Conclusion}

Kimi K2.5 is a large Mixture-of-Experts model that combines language understanding, multimodal capabilities, reasoning, coding, and agentic workflows.

Its approximately 1 trillion total parameters and approximately 32 billion active parameters demonstrate the use of sparse Mixture-of-Experts architecture.

Its 256K-token context window makes it suitable for long-context applications, while its native multimodal capabilities allow it to work with visual information.

One of its most interesting features is Agent Swarm, where complex tasks can be divided between multiple specialized agents.

Kimi K2.5 is therefore particularly relevant to modern AI applications involving:

\begin{itemize}
\item AI coding assistants
\item Multimodal AI
\item Research agents
\item Software engineering agents
\item Long-context applications
\item Tool-using AI systems
\end{itemize}

However, its large model size means that local deployment requires significant hardware. For most individual developers, API access is a more practical approach.

Overall, Kimi K2.5 demonstrates the current direction of advanced LLM development: moving from models that simply generate text toward systems that can understand multiple modalities, reason about complex problems, use tools, write code, and coordinate multiple agents.

\end{document}
